%
% 35-meromorph.tex -- Meromorphe Funktionen
%
% (c) 2026 Prof Dr Andreas Müller
%
\subsection{Meromorphe Funktionen}
Funktionen mit lauter isolierte Singulariäteten sind besonders
leicht zu behandeln und verdienen daher einen eigenen Namen.

\begin{definition}[meromorphe Funktion]
Eine holomorphe Funktion $f\colon U\to\mathbb{C}$,
für die $\mathbb{C}\setminus U$ aus lauter isolierten Punkten besteht,
heisst \emph{meromorph}.
\index{meromorph}%
\end{definition}

Die wichtigsten Beispiel für meromorphe Funktionen sind die rationalen
Funktionen
\[
f(z)
=
\frac{f(z)}{g(z)},
\]
wobei $f(z)$ und $g(z)$ Polynome sind.
Die Singularitäten sind dann genau die Nullstellen von $g(z)$ und die
Ordnung der Pole von $f(z)$ ist die Ordnung der Nullstellen von $g(z)$.


